\documentclass{article}
\usepackage{amsmath, amssymb, amsthm}
\newtheorem{theorem}{Theorem}
\newtheorem{lemma}[theorem]{Lemma}
\newtheorem{corollary}[theorem]{Corollary}
\newtheorem{proposition}[theorem]{Proposition}
\theoremstyle{definition}
\newtheorem{definition}[theorem]{Definition}
\begin{document}

\section*{The Fundamental Theorem of Finitely Generated Abelian Groups}

\subsection*{1.\ Definitions}

\begin{definition}[Group, abelian group]
A \emph{group} is a set $G$ with a binary operation
$\cdot:G\times G\to G$ that is associative, has a two-sided identity
$e$, and admits two-sided inverses. It is \emph{abelian} if
$a\cdot b=b\cdot a$ for all $a,b\in G$. Abelian groups will be written
additively, with identity $0$.
\end{definition}

\begin{definition}[Subgroup generated by a set]
For $S\subseteq G$, the \emph{subgroup generated by} $S$, written
$\langle S\rangle$, is the intersection of all subgroups of $G$ that
contain $S$. In the abelian case it consists of all finite
$\mathbb{Z}$-linear combinations $n_1s_1+\cdots+n_ks_k$ with
$n_i\in\mathbb{Z}$, $s_i\in S$.
\end{definition}

\begin{definition}[Finitely generated]
$G$ is \emph{finitely generated} if $G=\langle S\rangle$ for some
finite $S\subseteq G$.
\end{definition}

\begin{definition}[Cyclic group]
$G$ is \emph{cyclic} if $G=\langle a\rangle$ for some $a\in G$. Up to
isomorphism the cyclic groups are $\mathbb{Z}$ and $\mathbb{Z}/n\mathbb{Z}$ for $n\ge 1$.
\end{definition}

\begin{definition}[Direct sum]
The (external) \emph{direct sum} $A_1\oplus\cdots\oplus A_k$ is the
set of tuples $(a_1,\dots,a_k)$ with $a_i\in A_i$ under coordinatewise
addition.
\end{definition}

\begin{definition}[Free abelian group, basis, rank]
An abelian group $F$ is \emph{free of rank} $n$ if $F\cong\mathbb{Z}^n$.
A subset $\{e_1,\dots,e_n\}\subseteq F$ is a \emph{basis} if every
$x\in F$ has a unique expression $x=\sum c_ie_i$ with $c_i\in\mathbb{Z}$.
The rank is well-defined: $F\otimes_{\mathbb{Z}}\mathbb{Q}\cong\mathbb{Q}^n$ has
$\mathbb{Q}$-dimension $n$.
\end{definition}

\begin{definition}[Torsion]
$g\in G$ is \emph{torsion} if $ng=0$ for some $n\ge 1$. The torsion
elements form a subgroup $T(G)$, the \emph{torsion subgroup}; $G$ is
\emph{torsion-free} if $T(G)=0$.
\end{definition}

\subsection*{2.\ Statement}

\begin{theorem}[Fundamental Theorem]\label{thm:main}
Let $G$ be a finitely generated abelian group. There exist unique
integers $r\ge 0$, $k\ge 0$, and $n_1,\dots,n_k\ge 2$ with
$n_1\mid n_2\mid\cdots\mid n_k$ such that
\[
  G\;\cong\;\mathbb{Z}^r\oplus\mathbb{Z}/n_1\mathbb{Z}\oplus\mathbb{Z}/n_2\mathbb{Z}
  \oplus\cdots\oplus\mathbb{Z}/n_k\mathbb{Z}.
\]
The integer $r$ is the \emph{free rank}; the $n_i$ are the
\emph{invariant factors}.
\end{theorem}

\subsection*{3.\ Auxiliary lemmas}

\begin{lemma}\label{lem:fgfree}
Every subgroup $H\le\mathbb{Z}^n$ is free abelian of rank $\le n$.
\end{lemma}
\begin{proof}
Induction on $n$. For $n=1$, every subgroup of $\mathbb{Z}$ equals $d\mathbb{Z}$
for some $d\ge 0$ (take $d$ minimal positive, if any; division
algorithm), so it is free of rank $0$ or $1$. For $n>1$, let
$\pi:\mathbb{Z}^n\to\mathbb{Z}$ be the last-coordinate projection. Then
$\pi(H)=d\mathbb{Z}$. If $d=0$ then $H\subseteq\mathbb{Z}^{n-1}$ and induction
applies. Otherwise pick $h\in H$ with $\pi(h)=d$, and set
$K=H\cap\ker\pi\subseteq\mathbb{Z}^{n-1}$. By induction, $K$ has a basis
$e_1,\dots,e_m$. Every $x\in H$ satisfies $\pi(x)=qd$ for some
$q\in\mathbb{Z}$, and $x-qh\in K$ has a unique expansion in
$e_1,\dots,e_m$. Hence $e_1,\dots,e_m,h$ is a basis of $H$.
\end{proof}

\begin{lemma}\label{lem:pres}
Every finitely generated abelian group $G$ fits in a short exact
sequence $0\to K\to\mathbb{Z}^n\to G\to 0$ with $K$ free abelian of rank
$m\le n$.
\end{lemma}
\begin{proof}
If $g_1,\dots,g_n$ generate $G$, the map $\varphi:\mathbb{Z}^n\to G$,
$(a_i)\mapsto\sum a_ig_i$, is surjective; set $K=\ker\varphi$ and
apply Lemma~\ref{lem:fgfree}.
\end{proof}

\begin{definition}[Smith Normal Form]
$A\in M_{n\times m}(\mathbb{Z})$ is in \emph{Smith Normal Form} if it is
zero off the main diagonal and its diagonal entries
$d_1,\dots,d_s$ (with the remaining diagonal positions equal to $0$)
satisfy $d_1\mid d_2\mid\cdots\mid d_s$, all $d_i>0$.
\end{definition}

\begin{lemma}[Smith Normal Form]\label{lem:snf}
For every $A\in M_{n\times m}(\mathbb{Z})$ there exist $U\in GL_n(\mathbb{Z})$
and $V\in GL_m(\mathbb{Z})$ such that $UAV$ is in Smith Normal Form. The
diagonal entries $d_1,\dots,d_s$ are uniquely determined by $A$.
\end{lemma}

\begin{proof}
\textbf{Existence.} The integer elementary row/column operations
(row/column swaps; multiplication by $\pm 1$; adding an integer
multiple of one row/column to another) correspond to left/right
multiplication by elements of $GL_n(\mathbb{Z})$, $GL_m(\mathbb{Z})$.

If $A=0$ we are done. Otherwise let $\delta(A)$ be the smallest
positive absolute value among the entries of $A$, and move an entry
achieving it to position $(1,1)$, with sign chosen so that
$a_{11}=\delta(A)>0$. For each $j>1$, write $a_{1j}=qa_{11}+r$ with
$0\le r<a_{11}$ (division algorithm) and subtract $q$ times column $1$
from column $j$; if $r>0$ the new $a_{1j}=r$ achieves a strictly
smaller $\delta$, so restart. Do the same to clear column $1$ via row
operations. Since $\delta$ is a positive integer that strictly
decreases at each restart, the process terminates with row $1$ and
column $1$ entirely zero except $a_{11}>0$.

If some entry $a_{ij}$ ($i,j>1$) is not divisible by $a_{11}$, add
row $i$ to row $1$ and repeat the previous step; again $\delta$
strictly decreases. After finitely many iterations $a_{11}$ divides
all remaining entries; recursively apply the procedure to the
submatrix obtained by deleting row $1$ and column $1$. Since
divisibility by $a_{11}$ is preserved, the divisibility chain
$d_1\mid d_2\mid\cdots$ holds.

\textbf{Uniqueness.} For $1\le k\le\min(n,m)$ let $\Delta_k(A)$
denote the gcd of all $k\times k$ minors of $A$, with
$\Delta_0(A)=1$ and $\Delta_k(A)=0$ if every $k\times k$ minor
vanishes. Each elementary operation expresses the new $k\times k$
minors as integer linear combinations of the old ones (and
conversely), so $\Delta_k$ is invariant under multiplication by
elements of $GL_n(\mathbb{Z})$, $GL_m(\mathbb{Z})$. For a matrix in Smith form
with diagonal $d_1\mid\cdots\mid d_s$ one computes
$\Delta_k=d_1d_2\cdots d_k$ for $k\le s$ and $\Delta_k=0$ for $k>s$.
Hence $s$ is determined by $A$, and
$d_k=\Delta_k(A)/\Delta_{k-1}(A)$ is determined by $A$.
\end{proof}

\subsection*{4.\ Existence in Theorem~\ref{thm:main}}

By Lemma~\ref{lem:pres}, $G\cong\mathbb{Z}^n/K$ with $K\le\mathbb{Z}^n$ free of
rank $m\le n$. Choose a basis $f_1,\dots,f_m$ of $K$ and the
standard basis $e_1,\dots,e_n$ of $\mathbb{Z}^n$. Write
$f_j=\sum_iA_{ij}e_i$; this defines $A\in M_{n\times m}(\mathbb{Z})$.

By Lemma~\ref{lem:snf} there exist $U\in GL_n(\mathbb{Z})$,
$V\in GL_m(\mathbb{Z})$ with $UAV$ in Smith form, diagonal
$d_1\mid\cdots\mid d_s$, $s\le m$. The columns of $U^{-1}$ form a
new basis $e_1',\dots,e_n'$ of $\mathbb{Z}^n$, and the columns of $V$
describe a new basis $f_1',\dots,f_m'$ of $K$; in these bases
$f_j'=d_je_j'$ for $j\le s$. Since $K$ is free of rank $m$ we must
have $s=m$, and
\[
  K=\mathbb{Z}\,(d_1e_1')\oplus\cdots\oplus\mathbb{Z}\,(d_me_m')
  \;\subseteq\;\mathbb{Z} e_1'\oplus\cdots\oplus\mathbb{Z} e_n'=\mathbb{Z}^n,
\]
so
\[
  G\;\cong\;\mathbb{Z}^n/K\;\cong\;
  \mathbb{Z}/d_1\oplus\cdots\oplus\mathbb{Z}/d_m\oplus\mathbb{Z}^{\,n-m}.
\]
Discard the trivial summands $\mathbb{Z}/1=0$ and rename the remaining
$d_i\ge 2$ as $n_1\mid\cdots\mid n_k$. With $r=n-m$, this gives the
required decomposition. \qed

\subsection*{5.\ Uniqueness in Theorem~\ref{thm:main}}

\begin{lemma}\label{lem:rank}
If $G\cong\mathbb{Z}^r\oplus T$ with $T$ finite, then $T=T(G)$ and
$r=\dim_{\mathbb{Q}}(G\otimes_{\mathbb{Z}}\mathbb{Q})$. In particular $r$ depends
only on $G$.
\end{lemma}
\begin{proof}
Every element of $T$ has finite order; no nonzero element of
$\mathbb{Z}^r$ does. So $T(G)=T$. Tensoring with $\mathbb{Q}$ kills the
(finite) torsion and gives $\mathbb{Q}^r$; the dimension over $\mathbb{Q}$ is
an invariant of $G$.
\end{proof}

By Lemma~\ref{lem:rank}, $r$ is determined by $G$, and the torsion
part $T(G)\cong\mathbb{Z}/n_1\oplus\cdots\oplus\mathbb{Z}/n_k$ is determined as
a group. So we may assume $G$ is finite and prove uniqueness of the
invariant factors $n_1\mid\cdots\mid n_k$, all $n_i\ge 2$.

\begin{lemma}\label{lem:pdim}
Let $p$ be a prime and $i\ge 0$. The quotient
\[
  U_i(G)\;:=\;p^{\,i}G\,\big/\,p^{\,i+1}G
\]
is naturally an $\mathbb{F}_p$-vector space, and
\[
  \dim_{\mathbb{F}_p}U_i(G)
  \;=\;\#\{\,j:p^{\,i+1}\mid n_j\,\}.
\]
\end{lemma}
\begin{proof}
Multiplication by $p$ kills $U_i(G)$, giving the $\mathbb{F}_p$-structure.
Direct sums commute with $p^i(-)$ and quotients, so it suffices to
compute the contribution of a single summand $C=\mathbb{Z}/n$. We have
$p^iC=(p^i\mathbb{Z}+n\mathbb{Z})/n\mathbb{Z}=\gcd(p^i,n)\mathbb{Z}/n\mathbb{Z}$, which has order
$n/\gcd(p^i,n)$. Thus
\[
  |p^iC/p^{i+1}C|
  \;=\;\frac{|p^iC|}{|p^{i+1}C|}
  \;=\;\frac{\gcd(p^{i+1},n)}{\gcd(p^i,n)}.
\]
Write $n=p^{v}m$ with $\gcd(p,m)=1$. If $v\ge i+1$, the ratio equals
$p$; otherwise it equals $1$. So the summand $\mathbb{Z}/n_j$ contributes
$1$ to $\dim_{\mathbb{F}_p}U_i(G)$ exactly when $p^{i+1}\mid n_j$, proving
the formula.
\end{proof}

\begin{proof}[Proof of uniqueness]
Let $p$ be a prime and write $a_p(j)=v_p(n_j)$, the $p$-adic valuation
of $n_j$. Since $n_1\mid\cdots\mid n_k$, we have
$0\le a_p(1)\le a_p(2)\le\cdots\le a_p(k)$.
Lemma~\ref{lem:pdim} gives, for every $i\ge 0$,
\[
  \dim_{\mathbb{F}_p}U_i(G)\;=\;\#\{\,j:a_p(j)\ge i+1\,\}.
\]
The right-hand side, viewed as a function of $i$, completely determines
the nondecreasing sequence $(a_p(1),\dots,a_p(k))$. Indeed, the number
of $j$ with $a_p(j)\ge 1$ is $\dim U_0$, the number with
$a_p(j)\ge 2$ is $\dim U_1$, and so on; reading off these counts
recovers each $a_p(j)$. Doing this for every prime $p$ recovers each
$n_j=\prod_p p^{a_p(j)}$, and hence the entire list of invariant
factors. Since the $\dim_{\mathbb{F}_p}U_i(G)$ depend only on $G$ as an
abstract abelian group, the invariant factors are unique. \qed
\end{proof}

\subsection*{6.\ Conclusion}

Combining the existence and uniqueness arguments completes the proof of
Theorem~\ref{thm:main}. As a corollary, two finitely generated abelian
groups are isomorphic if and only if they have the same free rank $r$
and the same list of invariant factors $n_1\mid\cdots\mid n_k$.

\begin{corollary}[Primary decomposition]
Every finite abelian group is isomorphic to a direct sum of cyclic
groups of prime-power order; the multiset of prime-power orders is
unique.
\end{corollary}
\begin{proof}
Apply the Chinese Remainder Theorem to each $\mathbb{Z}/n_j$, writing
$n_j=\prod_pp^{a_p(j)}$, and regroup. Uniqueness follows from the
uniqueness of the exponents $a_p(j)$ established above.
\end{proof}

\end{document}